\answer{ex:18:1}{(a)~$0.9949\le w\le1.0049$, isto é, $|1-w|\lesssim0.005$. (b)~$K\ge400$ sinapses.}
\answer{ex:18:2}{$(\omega_R\tau_w)^2=\sqrt{\alpha(\alpha+2+2\varrho)}/\varrho-1$, $\varrho=\tau_m/\tau_w$; $f_R\approx11.1$~Hz; $|Z(\omega_R)|/|Z(0)|=4.6$.}
\answer{ex:18:3}{(a)~$f=1/\bigl(\tau\ln\frac{\mu}{\mu-\theta}\bigr)$; para $1$~Hz é preciso $\mu/\theta-1\approx3.7\cdot10^{-44}$. (b)~$T=\pi\tau/\sqrt\mu$, $f\approx3.2$~Hz: $f\propto\sqrt\mu$.}
\answer{ex:18:4}{O integrador funciona com $f\lesssim17.7$~Hz (filtro passa-baixas); o ressonador, só na faixa de $5.5$--$22$~Hz (filtro passa-faixa).}
\answer{ex:18:5}{(a)~$T_{\min}=\frac{\tau}{w-1}\ln\frac{0.5}{0.3}\approx10.2\ms$. (b)~$B>w-1-0.2=0.3$.}
\answer{ex:18:6}{Constante de tempo do ajuste $\tau/(\eta\langle r^2\rangle)$, $\eta=\tau\ln10/60\,\text{s}\approx3.8\cdot10^{-3}$; com $I\neq0$ o peso se desloca, e é preciso subtrair de $e$ uma cópia do comando $I/\tau$.}
\answer{ex:18:7}{Questão aberta. O ressonador ganha só $1.35$ vezes a $11$~Hz e perde $17$ vezes a $1$~Hz; o ganho principal da reconfiguração está na seleção de banda por limiar (exercício~\ref{ex:18:4}).}
